% !Mode:: "TeX::UTF-8"
\documentclass{svk_en}
%% ak pisete po slovensky, pouzijete namiesto horneho riadku
%% \documentclass{svk_sk}

\begin{document}
\title{Constructing flows from flow-pairs}

\author{Lukáš Gáborik
\email{gaborik5@uniba.sk}
\and 
Jozef Rajník
\email{rajnik@dcs.fmph.uniba.sk}
\thanks{Supervisor}}
%% vsimnite si, ze u autorov sa nepisu tituly
%% prikaz \inst sluzi ako odkaz do zoznamu institucii
%% (vid. nizsie)

%% ak je skolitel/ka prace sucastne autorom, oznacte ho/ju uvedenym sposobom
%% ak praca ma skolitela, ktory ale nie je autorom, uvedte to cez nasledujuci prikaz
%% \supervisor{Tomáš Vinař \inst{1}\email{vinar@fmph.uniba.sk}}
%% ak praca nema skolitela, jednoducho vynechajte


%% nasleduje kratka verzia nazvu clanku a 
%% zoznam autorov (bez krstnych mien)
%% tieto informacie sa zobrazuju v hlavicke
\titlerunning{Flow-pairs}
\authorrunning{Gáborik and Rajník}

\institute{Department of Computer Science,
FMFI UK,
Mlynská dolina,
842~48~Bratislava}

\maketitle

\begin{abstract}
A nowhere-zero $k$-flow is a mapping of values $\pm 1, \pm 2, \dots, \pm(k-1)$ to edges satisfying Kirchhoff's law -- the inflow is the same as the outflow -- for every vertex of the graph. It is conjectured \cite{tutte} that each bridgeless graph has a nowhere-zero $5$-flow. Almost three decades later, Seymour \cite{seymour} made a general construction of nowhere-zero $6$-flow from $2$-flow and $3$-flow. These two flows permit zeroes, but no edge can have both of them equal to zero. In a recent preprint \cite{chebyshev_flows_preprint}, along the introduction of two-dimensional Chebyshev nowhere-zero flows, the authors showed a similar construction of nowhere-zero $5$-flow from $2$-flow and $4$-flow. Except two simultaneous zeroes, flow values $(0, 1)$ and $(0, -1)$ are now forbidden, too. This structure is called a $1/2$-flow-pair.

In this paper, we explore the properties of $1/2$-flow-pairs. We prove their existence on snarks with oddness $2$, and we also show that one can construct a rich nowhere-zero $7$-flow from them. We also investigate the Chebyshev flow number and show that the two-dimensional Chebyshev flow number of snarks can be arbitrarily small.
\keywords{...}
\end{abstract}

\section*{Acknowledgments}
...

\nocite{*}
\bibliographystyle{apalike}
\bibliography{references}

%% citacie ulozte do suboru references.bib
%% na populaciu zoznamu literatury pouzite program
%%
%% bibtex latex_long_en
%% predosly prikaz samozrejme predpoklada,
%% ze sa subor ktory  "latex-ujete" vola latex_long_en
%%
%% po ktorom je potrebne dokument znova zlatexovat

\end{document}
